% \section{Conclusion}

% This paper presented a multilingual multimodal framework for cybercrime detection tailored to the Indian context. By integrating textual and visual modalities through a transformer-based fusion mechanism, the proposed approach effectively captures complex patterns across diverse data sources and languages. 

% Experimental results demonstrate that the model outperforms strong baseline methods, highlighting the benefits of multimodal learning and cross-lingual representation. The inclusion of explainable artificial intelligence techniques further enhances the transparency and practical applicability of the system.

% The proposed dataset and hierarchical taxonomy contribute valuable resources for future research in this domain. Overall, this work provides a scalable and robust solution for real-world cybercrime detection. Future work will focus on expanding the dataset, improving low-resource language performance, and exploring real-time deployment scenarios.

\section{Conclusion}
This paper presented a transformer-based multilingual 
multimodal framework for cybercrime detection in the 
Indian context, combining Qwen2.5-7B text encoding, 
ViT-B/16 visual encoding, and cross-modal attention 
fusion across a novel 25,000-sample dataset spanning 
five languages and 50 crime categories. Experimental 
results demonstrate a macro F1 of 0.883, outperforming
all baselines by up to 5.3 points, with LIME and
attention maps providing interpretability. Future work
will focus on data
augmentation for low-resource languages, real-time 
deployment, and expanding the dataset to additional 
Indian languages.